\chapter{Análisis del Valor y Aplicabilidad Industrial} \label{ch:valor-aplicabilidad}

\section{Valor para redes OT reales} \label{sec:valor-redes-reales}

El valor principal de GEMEROTIC se encuentra en reducir el riesgo de actuar sobre redes OT sin una representación técnica suficiente. En organizaciones industriales con varias plantas, activos heterogéneos y continuidad productiva crítica, el inventario incompleto y la falta de trazabilidad dificultan responder preguntas básicas: qué equipos existen, cómo están conectados, qué zonas comunican entre sí, qué cambios se han ensayado y qué evidencias respaldan una decisión.

Un gemelo digital no elimina la necesidad de auditoría, pero puede reducir el coste y el riesgo de prepararla. Permite modelar la red, reproducir escenarios, revisar cambios y conservar evidencias sin interrumpir la operación real. Para una organización como AMC Global, este enfoque puede apoyar procesos de ciber-resiliencia, mejora continua y preparación ante marcos regulatorios.

\section{Comparación frente a enfoques tradicionales} \label{sec:comparacion-enfoques}

Frente a inventarios manuales o diagramas estáticos, GEMEROTIC propone un modelo persistente y transformable. Frente a una automatización directa, prioriza una fase previa de representación, validación y documentación. Esta diferencia es clave en OT: el objetivo inicial no es cambiar la red más rápido, sino entenderla mejor y reducir el riesgo de cualquier intervención.

El enfoque también se diferencia de una simulación aislada. El gemelo no se limita a levantar nodos; conecta inventario, topología, seguridad OT, artefactos y evidencias. Esa integración permite que el modelo sirva tanto para ingeniería de red como para ciberseguridad y auditoría técnica.

\section{Aplicabilidad a sectores críticos} \label{sec:sectores-criticos}

Aunque el caso motivador se sitúa en el sector alimentario y de bebidas, el enfoque es aplicable a otros sectores con redes OT: energía, agua, manufactura, transporte, logística, farmacéutica o infraestructuras públicas. En todos ellos aparecen restricciones similares: procesos continuos, ventanas de mantenimiento limitadas, activos legacy, exigencias regulatorias y necesidad de justificar decisiones técnicas.

La arquitectura modular permite adaptar el catálogo de activos, los perfiles runtime y los controles de cumplimiento a cada sector. En una planta alimentaria pueden priorizarse líneas de producción, SCADA, PLC y redes de planta; en energía pueden aparecer subestaciones, relés, telecontrol y protocolos específicos; en agua pueden destacarse estaciones remotas, sensores y telemetría.

\section{Encaje con estándares y auditorías} \label{sec:encaje-estandares}

GEMEROTIC no pretende certificar automáticamente el cumplimiento de ISA/IEC 62443, NIS2, NIST SP 800-82 o ENS. Su aportación consiste en estructurar evidencias técnicas que ayuden a una organización a prepararse para auditorías, revisiones internas y procesos de mejora continua.

El modelo de zonas y conductos facilita explicar segmentación. El inventario derivado ayuda a documentar activos. La persistencia granular permite conservar revisiones. Los artefactos reproducibles permiten justificar qué laboratorio o configuración se ha evaluado. El compliance determinista aporta hallazgos trazables. En conjunto, la plataforma puede convertirse en una herramienta de apoyo para responsables de seguridad, ingeniería y operación.

\section{Evolución empresarial} \label{sec:evolucion-empresarial}

Desde una perspectiva de producto, GEMEROTIC puede evolucionar hacia una plataforma de documentación, validación y asistencia a la auditoría para redes OT. Su valor diferencial reside en combinar un constructor visual, un modelo formal, inventario enriquecido, despliegue de laboratorio, validación consultiva y semántica normativa industrial.

La colaboración con una empresa industrial permite orientar la evolución hacia necesidades reales: integración con inventarios existentes, plantillas de planta, capturas de evidencias, perfiles por tipo de activo, validación de cambios y cuadros de mando de ciber-resiliencia. Esta orientación evita desarrollar una herramienta técnicamente interesante pero desconectada del uso operativo.
